\chapter{Cho và Nhận}
\markboth{Cho và Nhận}{Giving and Receiving}

\section{Một người cho rất nhiều nhưng lại cho để được biết ơn.}
\begin{truyen}{Một người cho rất nhiều nhưng lại cho để được biết ơn.}{Giving can hide a demand.}
\chuhoa{B}{ề ngoài, cô luôn giúp mọi người. Nhưng mỗi lần không được cảm ơn như ý, cô lại hụt hẫng và giận dỗi. Dần dần cô nhận ra mình không chỉ cho đi; mình còn âm thầm đòi lại bằng sự ghi nhận.}
\textit{(On the surface, she always helped people. Yet whenever she was not thanked as expected, she felt wounded and resentful. Slowly she realized that she had not been simply giving; she had been quietly asking to be repaid with recognition.)}
\end{truyen}
\begin{baihoc}
Cho đi mà buộc người khác mắc nợ cảm xúc thì chưa hẳn là rộng lượng.
\textit{(Giving while forcing others into emotional debt is not yet generosity.)}
\end{baihoc}
\begin{ghinhoanh}
\textbf{Short English to remember:}

Giving can carry hidden demands.
\end{ghinhoanh}
\begin{tuhocnhanh}
1. Khi giúp ai, em có mong họ đáp lại theo một cách cụ thể không? \textit{(When you help someone, do you expect a specific response?)}

2. Em đang cho thật hay đang trao đổi mà không gọi đúng tên? \textit{(Are you truly giving, or are you exchanging without naming it?)}
\end{tuhocnhanh}
\ngancach

\section{Một học sinh biết nhận lời góp ý mà không tự ái.}
\begin{truyen}{Một học sinh biết nhận lời góp ý mà không tự ái.}{Receiving correction is also a strength.}
\chuhoa{E}{m bị thầy sửa bài trước lớp. Ban đầu em đỏ mặt, nhưng thay vì chống đỡ, em ghi lại từng lỗi và sửa dần. Chính khả năng nhận điều khó nghe đã giúp em tiến nhanh hơn nhiều bạn chỉ thích được khen.}
\textit{(The student was corrected publicly. At first there was embarrassment, but instead of defending, the errors were written down and repaired. The ability to receive difficult words helped this student progress faster than many classmates who only liked praise.)}
\end{truyen}
\begin{baihoc}
Biết nhận cũng quý như biết cho.
\textit{(Knowing how to receive is as valuable as knowing how to give.)}
\end{baihoc}
\begin{ghinhoanh}
\textbf{Short English to remember:}

Receiving truth is a skill.
\end{ghinhoanh}
\begin{tuhocnhanh}
1. Em phản ứng thế nào khi bị góp ý? \textit{(How do you react when corrected?)}

2. Em cần học nhận điều gì với thái độ trưởng thành hơn? \textit{(What do you need to learn to receive more maturely?)}
\end{tuhocnhanh}
\ngancach

\section{Một người cho thời gian quý hơn cho tiền.}
\begin{truyen}{Một người cho thời gian quý hơn cho tiền.}{Time is often the deeper gift.}
\chuhoa{K}{hi bạn thất tình, nhiều người nhắn vài câu an ủi rồi thôi. Chỉ có một người ngồi cạnh suốt buổi chiều, không vội khuyên, không giảng đạo lý. Đôi khi điều cứu một người không phải là món quà lớn, mà là việc có ai đó dám cho họ thời gian thật sự.}
\textit{(When a friend was heartbroken, many sent a few comforting messages and moved on. Only one person stayed through the afternoon, not rushing to advise or preach. Sometimes what saves a person is not a big gift, but someone willing to give real time.)}
\end{truyen}
\begin{baihoc}
Thứ được cho bằng giờ sống thường chạm sâu hơn thứ được cho bằng tiền.
\textit{(What is given through living time often reaches deeper than what is given through money.)}
\end{baihoc}
\begin{ghinhoanh}
\textbf{Short English to remember:}

Time can be the purest gift.
\end{ghinhoanh}
\begin{tuhocnhanh}
1. Gần đây em đã cho ai thời gian thật sự chưa? \textit{(Have you recently given anyone your real time?)}

2. Khi em mệt, em cần lời khuyên hay cần một người ở cạnh? \textit{(When you are tired, do you need advice or presence?)}
\end{tuhocnhanh}
\ngancach

\section{Nhận sự giúp đỡ mà vẫn giữ lòng biết ơn.}
\begin{truyen}{Nhận sự giúp đỡ mà vẫn giữ lòng biết ơn.}{To receive well is to remember well.}
\chuhoa{C}{ó người được nâng đỡ lúc khó khăn rồi khi ổn hơn lại cư xử như thể mọi thứ là đương nhiên. Cũng có người chỉ nhận một lần giúp nhỏ nhưng nhớ rất lâu, rồi sống tử tế với người khác như một cách nối tiếp điều tốt đã nhận. Chính lòng biết ơn quyết định việc nhận làm ta lớn lên hay hư đi.}
\textit{(Some people are helped in hardship and later act as though it was all owed to them. Others receive one small act of support, remember it for years, and then treat others kindly as a continuation of what they received. Gratitude determines whether receiving makes us grow or decay.)}
\end{truyen}
\begin{baihoc}
Nhận mà không biết ơn dễ làm con người cạn dần phẩm chất.
\textit{(Receiving without gratitude easily drains character.)}
\end{baihoc}
\begin{ghinhoanh}
\textbf{Short English to remember:}

Gratitude completes receiving.
\end{ghinhoanh}
\begin{tuhocnhanh}
1. Em còn nhớ ai từng đỡ mình ở giai đoạn khó không? \textit{(Do you still remember who supported you in a difficult season?)}

2. Em đã nối tiếp điều tốt mình nhận được bằng cách nào? \textit{(How have you passed on the good you received?)}
\end{tuhocnhanh}
\ngancach

\section{Cho đúng cách để người nhận không thấy bị hạ thấp.}
\begin{truyen}{Cho đúng cách để người nhận không thấy bị hạ thấp.}{How you give matters as much as what you give.}
\chuhoa{M}{ột thầy giáo hỗ trợ học phí cho học sinh khó khăn nhưng không bao giờ nhắc lại trước lớp. Nhờ vậy, sự giúp đỡ đi vào như một bàn tay đỡ, không phải như một cái loa phóng to sự thiếu thốn của người nhận. Cho đúng cách là vừa giúp người, vừa giữ phẩm giá cho họ.}
\textit{(A teacher helped a struggling student with school fees but never mentioned it in public. In that way, the help entered like a supporting hand, not a loudspeaker magnifying the student's poverty. To give well is to help while preserving dignity.)}
\end{truyen}
\begin{baihoc}
Cho đi đẹp nhất là khi người nhận thấy được nâng lên chứ không bị đặt thấp xuống.
\textit{(Giving is most beautiful when the receiver feels lifted, not lowered.)}
\end{baihoc}
\begin{ghinhoanh}
\textbf{Short English to remember:}

A good gift protects dignity.
\end{ghinhoanh}
\begin{tuhocnhanh}
1. Khi giúp ai, em có vô tình làm họ ngại không? \textit{(When helping someone, do you unintentionally make them feel ashamed?)}

2. Làm sao để lòng tốt của em đi cùng sự tinh tế? \textit{(How can your kindness be joined with delicacy?)}
\end{tuhocnhanh}
\ngancach
